% Generated by scripts/24_program_overview.py; do not edit.
\begin{longtable}{@{}>{\raggedright\arraybackslash}p{1.3cm}>{\raggedright\arraybackslash}p{1.3cm}>{\raggedright\arraybackslash}p{1.4cm}>{\raggedright\arraybackslash}p{5.5cm}>{\raggedright\arraybackslash}p{1.3cm}>{\raggedright\arraybackslash}p{2.0cm}>{\raggedright\arraybackslash}p{1.7cm}@{}}
\caption{Signalling readouts and negative-control pathways (\protect\nolinkurl{data/curated/signalling_evidence.tsv}): role, type of evidence and its strength (strong: zebrafish, direct or loss/gain of the pathway; moderate: indirect zebrafish or other vertebrate evidence). Readouts mark cells that receive a signal; ligands and secreted antagonists mark cells that send or block one. Profile: the curated signalling profile of the gene (or why it is not used). Negative controls are pathways that should be inactive or uniform at 3--12\,hpf, and dissociation-induced genes.}\label{tab:signalling-evidence}\\
\toprule
Pathway & Gene & Role & Evidence & Strength & Profile & References \\
\midrule
\endfirsthead
\toprule
Pathway & Gene & Role & Evidence & Strength & Profile & References \\
\midrule
\endhead
\bottomrule
\endlastfoot
\multicolumn{7}{@{}l}{\textit{readouts (receiving)}} \\
BMP & \emph{id1} & target & direct target (Smad-binding BMP response element in the Id1 promoter; induced by BMP in ES cells), mouse only & moderate & BMP & \cite{hollnagel1999,korchynskyi2002} \\
\rowcolor{gray!10}BMP & \emph{id3} & target & direct target (induced by BMP2/7 protein under cycloheximide; reduced in bmp7 mutants) & strong & BMP & \cite{greenfeld2021,hollnagel1999} \\
BMP & \emph{bambia} & feedback & lost in swirl/bmp2b and expanded in chordino mutants; direct target (cycloheximide-resistant BMP induction) & strong & BMP & \cite{tsang2000,greenfeld2021} \\
\rowcolor{gray!10}BMP & \emph{szl} & feedback & reduced or absent in swirl/bmp2b and expanded in chordino mutants; direct target & strong & BMP & \cite{yabe2003,greenfeld2021} \\
BMP & \emph{msx1b} & target & Msx expression changes with BMP level (BMP mutants, dominant-negative receptor); induced at intermediate BMP, so not monotonic & moderate & BMP & \cite{tribulo2003,phillips2006} \\
\rowcolor{gray!10}BMP & \emph{smad6a} & feedback & direct target in zebrafish (cycloheximide-resistant BMP induction); BMP-responsive element in the mouse Smad6 promoter & strong & BMP & \cite{greenfeld2021,ishida2000} \\
BMP & \emph{smad6b} & feedback & direct target in zebrafish (cycloheximide-resistant BMP induction); BMP-responsive element in the mouse Smad6 promoter & strong & BMP & \cite{greenfeld2021,ishida2000} \\
\rowcolor{gray!10}BMP & \emph{ved} & target & direct BMP target (cycloheximide); also downstream of Wnt3a/Wnt8 (lost in wnt3a/wnt8 morphants) & strong & not used: multi-readout and Wnt8 target & \cite{greenfeld2021,shimizu2002,shimizu2005,ramel2005} \\
BMP & \emph{vox} & target & direct BMP target (cycloheximide); BMP-dependent and BMP-independent domains in BMP mutants; direct Wnt8/beta-catenin target & strong & not used: multi-readout and Wnt8 target & \cite{greenfeld2021,melby2000,ramel2004,ramel2005,imai2001} \\
\rowcolor{gray!10}BMP & \emph{vent} & target & direct BMP target (cycloheximide); mostly BMP-dependent in BMP mutants; direct Wnt8/beta-catenin target & strong & not used: multi-readout and Wnt8 target & \cite{greenfeld2021,melby2000,ramel2004,ramel2005,imai2001} \\
Nodal & \emph{lft1} & feedback & expression depends on Nodal signalling (antivin = lefty1); maintained by Cyclops; Lefty knockdown releases Nodal & strong & Nodal & \cite{meno1999,feldman2002,bisgrove1999,thisse1999} \\
\rowcolor{gray!10}Nodal & \emph{lft2} & feedback & Nodal-dependent transcription (Cyclops required for maintenance); Lefty knockdown releases Nodal & strong & Nodal & \cite{feldman2002,bisgrove1999} \\
Nodal & \emph{ndr1} & auto-regulation & positive autoregulation through a Nodal response element (FoxH1/Smad) in the squint first intron (transgenic reporters) & strong & Nodal & \cite{fan2007,feldman2002,rebagliati1998} \\
\rowcolor{gray!10}Nodal & \emph{ndr2} & auto-regulation & Nodal ligands induce their own transcription (positive autoregulatory loop), stated for both zebrafish Nodals & moderate & Nodal & \cite{feldman2002,rebagliati1998} \\
Wnt & \emph{sp5l} & target & direct Wnt/beta-catenin target (microarray, induction under cycloheximide); required downstream of wnt8 & strong & Wnt & \cite{weidinger2005,thorpe2005} \\
\rowcolor{gray!10}Wnt & \emph{sp5a} & target & acts downstream of Wnt8 with sp5l (combined knockdown phenocopies loss of wnt8) & moderate & Wnt & \cite{weidinger2005} \\
Wnt & \emph{axin2} & feedback & direct target (Tcf/Lef sites in the promoter and intron; rapid induction by Wnt), mouse & moderate & Wnt & \cite{jho2002} \\
\rowcolor{gray!10}Wnt & \emph{nkd1} & feedback & expression mirrors Wnt/beta-catenin activity (ventrolateral margin, axial mesendoderm); Nkd1 inhibits canonical Wnt & moderate & not used: multi-readout and BMP target & \cite{vanraay2007,greenfeld2021} \\
Wnt & \emph{notum1a} & feedback & Wnt/beta-catenin-induced secreted inhibitor (self-regulatory loop; Wnt gain and loss in zebrafish) & strong & Wnt & \cite{flowers2012} \\
\rowcolor{gray!10}Wnt & \emph{lef1} & target & direct Wnt target through a Wnt-responsive element in the LEF-1 promoter (human cells); in zebrafish Lef1 drives the TOPdGFP Wnt reporter & moderate & Wnt & \cite{filali2002,dorsky2002} \\
FGF & \emph{spry2} & feedback & induced by FGF (gain and loss of FGF signalling), mouse and chick & moderate & FGF & \cite{minowada1999} \\
\rowcolor{gray!10}FGF & \emph{spry4} & feedback & FGF feedback-induced antagonist: lost in ace/fgf8 mutants and with SU5402, induced by FGF8 & strong & FGF & \cite{furthauer2001} \\
FGF & \emph{spry1} & feedback & induced by FGF (vertebrate Sprouty family), mouse and chick & moderate & FGF & \cite{minowada1999} \\
\rowcolor{gray!10}FGF & \emph{dusp6} & feedback & FGF synexpression; under FGF control after initiation by maternal beta-catenin on the dorsal side; FGF reporter transgene & strong & FGF & \cite{tsang2004,molina2007} \\
FGF & \emph{etv4} & target & lost in fgf8 mutants and with SU5402; induced by ectopic fgf8 & strong & FGF & \cite{roehl2001,raible2001} \\
\rowcolor{gray!10}FGF & \emph{etv5b} & target & lost in fgf8 mutants and with SU5402; induced by ectopic fgf8 & strong & FGF & \cite{roehl2001,raible2001} \\
FGF & \emph{il17rd} & feedback & FGF synexpression; positively regulated by FGF; feedback inhibitor of Ras/MAPK & strong & FGF & \cite{tsang2002,furthauer2002} \\
\rowcolor{gray!10}RA & \emph{cyp26a1} & feedback & induced by exogenous RA (zebrafish differential display; time- and dose-dependent up-regulation) & moderate & RA (blocking) & \cite{white1996,dobbsmcauliffe2004,kudoh2002} \\
RA & \emph{dhrs3a} & feedback & RA-dependent and strongly RA-inducible (microarray: opposite response to RA and RA antagonist; in situ) & strong & RA (receiving) & \cite{feng2010} \\
\rowcolor{gray!10}Notch & \emph{her4.1} & target & activated by Notch1 intracellular domain (RNA injection); her4 regulatory elements report Notch activity (transgenic) & strong & Notch (receiving) & \cite{takke1999,yeo2007} \\
Notch & \emph{her4.2} & target & as HER4.1 (her4 cluster) & strong & Notch (receiving) & \cite{takke1999,yeo2007} \\
\rowcolor{gray!10}Notch & \emph{her4.3} & target & as HER4.1 (her4 cluster) & strong & Notch (receiving) & \cite{takke1999,yeo2007} \\
Notch & \emph{her4.4} & target & as HER4.1 (her4 cluster) & strong & Notch (receiving) & \cite{takke1999,yeo2007} \\
\rowcolor{gray!10}Notch & \emph{her15.1} & target & PSM expression regulated by Delta-Notch (mutants, morpholinos, inhibitor), in a complex manner; depends on her7 & moderate & Notch (receiving) & \cite{shankaran2007} \\
Notch & \emph{her15.2} & target & as HER15.1 & moderate & Notch (receiving) & \cite{shankaran2007} \\
\rowcolor{gray!10}Notch & \emph{her12} & target & PSM expression regulated by Delta-Notch (mutants, morpholinos, inhibitor), complex & moderate & Notch (receiving) & \cite{shankaran2007} \\
Notch & \emph{nrarpa} & feedback & direct Notch target (Xenopus: Su(H)-dependent induction); zebrafish nrarpa also stabilises Lef1 (Wnt) & moderate & not used: multi-readout, also acts in Wnt/Lef1 & \cite{lamar2001,ishitani2005} \\
\rowcolor{gray!10}Notch & \emph{nrarpb} & feedback & direct Notch target (Xenopus); direct BMP target in zebrafish (cycloheximide) & moderate & not used: multi-readout and BMP target & \cite{lamar2001,greenfeld2021} \\
Notch & \emph{hey1} & target & dynamic PSM expression impaired in deltaD (after-eight) mutants; ectopic deltaD induces hey1; promoter is a direct Notch target in mouse and human & strong & Notch (receiving) & \cite{winkler2003,maier2000} \\
\rowcolor{gray!10}Hedgehog & \emph{ptch1} & feedback & transcription depends on Hh signalling; responds to lower Hh levels than ptch2 & strong & Hedgehog (receiving) & \cite{lewis1999} \\
Hedgehog & \emph{ptch2} & feedback & transcription follows shh activity; ectopic shh activates it; PKA-dependent & strong & Hedgehog (receiving) & \cite{concordet1996,lewis1999} \\
\rowcolor{gray!10}Hedgehog & \emph{gli1} & target & Shh activates Gli1 transcription (Xenopus); in zebrafish Gli1 (detour) is an activator of Hh target genes & moderate & Hedgehog (receiving) & \cite{lee1997,karlstrom2003} \\
Hedgehog & \emph{hhip} & feedback & induced by ectopic Hh and lost in Hh mutants (mouse); zebrafish hhip is expressed in Hh-responding adaxial cells & moderate & Hedgehog (receiving) & \cite{chuang1999hhip,ochi2006} \\
\midrule
\multicolumn{7}{@{}l}{\textit{negative controls}} \\
IFN & \emph{isg15} & control & IFN-stimulated gene (induced by recombinant type I IFN and by virus via IFNR in larvae); also induced by morpholino injection & moderate & -- & \cite{levraud2019,lai2019,lam2004,herbomel1999} \\
\rowcolor{gray!10}IFN & \emph{mxa} & control & IFN-stimulated gene (larvae, recombinant IFN and virus) & moderate & -- & \cite{levraud2019,lam2004,herbomel1999} \\
IFN & \emph{rsad2} & control & IFN-stimulated gene (viperin; larvae, recombinant IFN and virus) & moderate & -- & \cite{levraud2019,lam2004,herbomel1999} \\
\rowcolor{gray!10}AhR & \emph{cyp1a} & control & AHR target induced by TCDD and PCB126 in embryos (xenobiotic response elements) & strong & -- & \cite{andreasen2002,jonsson2007} \\
AhR & \emph{cyp1b1} & control & induced by PCB126 in embryos (xenobiotic response elements) & moderate & -- & \cite{jonsson2007,andreasen2002} \\
\rowcolor{gray!10}AhR & \emph{cyp1c1} & control & induced by PCB126 in embryos; high basal embryonic expression & moderate & -- & \cite{jonsson2007} \\
AhR & \emph{ahrrb} & control & induced by TCDD in ZF-L cells and embryos; AHR negative-feedback repressor & strong & -- & \cite{evans2005} \\
\rowcolor{gray!10}GR & \emph{fkbp5} & control & GR target: up-regulated by full-length GR mRNA at 6 hpf; GR-dependent and cortisol-induced in larvae & strong & -- & \cite{gans2020,gans2021,alsop2008} \\
GR & \emph{klf9} & control & under-expressed in GR mutant larvae and induced by cortisol; feedforward regulator of fkbp5 & strong & -- & \cite{gans2020,gans2021,alsop2008} \\
\rowcolor{gray!10}GR & \emph{tsc22d3} & control & described as a GR target (gilz) with GR-dependent daily oscillation in larvae & moderate & -- & \cite{gans2021,alsop2008} \\
circadian & \emph{per1b} & control & zygotic period1 transcription requires CLOCK (dominant-negative CLOCK blocks it) & moderate & -- & \cite{dekens2008} \\
\rowcolor{gray!10}circadian & \emph{per2} & control & light-induced in all zebrafish cells through a D-box/E-box light-responsive module; light at 0-6 hpf transiently raises per2 & moderate & -- & \cite{vatine2009,ziv2006} \\
circadian & \emph{cry1aa} & control & light-induced clock repressor (cry1a); induction correlates with clock resetting & moderate & -- & \cite{tamai2007} \\
\rowcolor{gray!10}circadian & \emph{tefa} & control & light-inducible D-box factor that drives per2 light induction & moderate & -- & \cite{vatine2009} \\
dissoc. & \emph{fosab} & control & immediate-early gene induced by enzymatic dissociation at 37 C (scRNA-seq), mouse & moderate & -- & \cite{vandenbrink2017,denisenko2020} \\
\rowcolor{gray!10}dissoc. & \emph{junba} & control & immediate-early gene induced by dissociation (mouse) & moderate & -- & \cite{vandenbrink2017,denisenko2020} \\
dissoc. & \emph{jun} & control & immediate-early gene induced by dissociation (mouse) & moderate & -- & \cite{vandenbrink2017,denisenko2020} \\
\rowcolor{gray!10}dissoc. & \emph{atf3} & control & immediate-early gene induced by dissociation (mouse) & moderate & -- & \cite{vandenbrink2017,denisenko2020} \\
dissoc. & \emph{egr1} & control & immediate-early gene induced by dissociation (mouse) & moderate & -- & \cite{vandenbrink2017,denisenko2020} \\
\midrule
\multicolumn{7}{@{}l}{\textit{ligands (sending)}} \\
Nodal & \emph{ndr1} & ligand & in situ: maternal, then blastoderm margin at blastula; dorsal blastomeres at MBT and dorsal YSL, whole margin and YSL in early gastrula, dorsal forerunner cells at shield; sqt;cyc double mutants lack organizer, mesoderm and endoderm; misexpression (also only in the YSL) expands organizer fates; Squint/Cyclops required in the YSL for endoderm and head mesoderm & strong & Nodal & \cite{rebagliati1998,erter1998,feldman1998,fan2007} \\
\rowcolor{gray!10}Nodal & \emph{ndr2} & ligand & in situ: blastoderm margin at blastula; dorsal mesendoderm and shield hypoblast at gastrula, prechordal plate until early somitogenesis, later left lateral plate and left diencephalon; cyc mutants lack floor plate and ventral brain, rescued non-autonomously by cyc RNA; redundant with sqt for organizer and mesendoderm & strong & Nodal & \cite{rebagliati1998,sampath1998,feldman1998} \\
BMP & \emph{bmp2b} & ligand & in situ: whole blastula, restricted ventrally before gastrulation (Fgf-dependent); ventral gastrula; swirl/bmp2b mutants dorsalized (loss of non-neural ectoderm, neural crest, ventral mesoderm), cell non-autonomous; overexpression ventralizes & strong & BMP & \cite{kishimoto1997,nguyen1998,nikaido1997,furthauer2004,schmid2000,pyati2005} \\
\rowcolor{gray!10}BMP & \emph{bmp4} & ligand & in situ: ventral gastrula and also the embryonic shield; null mutants affect late-gastrula specification of ventroposterior fates & strong & BMP & \cite{nikaido1997,stickney2007} \\
BMP & \emph{bmp7a} & ligand & snailhouse/bmp7 mutants strongly dorsalized, null identical to swirl/bmp2b; Bmp2b/Bmp7 synergize (likely heterodimers), transduced by Smad5 and antagonized by Chordino; Bmp genes expressed in the whole blastula, then restricted ventrally & strong & BMP & \cite{dick2000,schmid2000,furthauer2004} \\
\rowcolor{gray!10}BMP & \emph{admp} & ligand & in situ: dorsal only, from blastula in the organizer region (anterior neuroectoderm and axial mesoderm precursors), vanishing in an anterior-to-posterior wave during gastrulation; also tail organizer; knockdown enlarges the organizer and impairs ventral/posterior development, overexpression is Bmp-like; cooperates with Bmp2b/Bmp7 & strong & BMP organizer (sending) & \cite{lele2001,willot2002} \\
Wnt & \emph{wnt8a} & ligand & in situ: maternal and zygotic; blastoderm margin at 50\% epiboly (cleared from the dorsal midline in late gastrula), tailbud during segmentation; deficiency and morpholino: DV mesoderm and AP neural patterning defects; wnt3a/wnt8 double knockdown expands the organizer and removes the posterior body & strong & Wnt & \cite{kelly1995,lekven2001,ramel2005,shimizu2005} \\
\rowcolor{gray!10}Wnt & \emph{wnt3a} & ligand & in situ: blastoderm margin (overlapping wnt8) and later tailbud; wnt3a/wnt8 double knockdown anteriorizes the neuroectoderm, expands the organizer and blocks tail formation; sp5l in the tail depends on wnt3a and wnt8 & strong & Wnt & \cite{shimizu2005,thorpe2005} \\
Wnt/PCP & \emph{wnt5b} & ligand & maternal and zygotic at all gastrula stages; pipetail mutants: cell elongation and convergent extension defects in posterior regions of the gastrula (redundant with wnt11 anteriorly); Frizzled-2 acts downstream & strong & Wnt/PCP (sending) & \cite{kilian2003} \\
\rowcolor{gray!10}Wnt/PCP & \emph{wnt11} & ligand & in situ: blastoderm margin (presumptive mesoderm) from late blastula, induced by the yolk cell and FGF, independent of ntl; later somites and notochord; silberblick mutants: convergent extension defects and cyclopia, cell non-autonomous, rescued by a Dishevelled that does not signal through the canonical pathway & strong & Wnt/PCP (sending) & \cite{makita1998,heisenberg2000} \\
FGF & \emph{fgf8a} & ligand & in situ: dorsoventral gradient (dorsal high) at the gastrula margin, not in axial mesoderm; mesodermal precursors; MHB and somitogenesis domains; overexpression dorsalizes and represses ventral bmp2b/bmp4; acerebellar mutants lose the MHB organizer, with mild DV and mesoderm defects; with fgf24 required for posterior mesoderm & strong & FGF & \cite{furthauer1997,reifers1998,draper2003,furthauer2004} \\
\rowcolor{gray!10}FGF & \emph{fgf24} & ligand & in situ: co-expressed with fgf8 in mesodermal precursors during gastrulation; fgf24 knockdown in fgf8 mutants blocks most posterior mesoderm & strong & FGF & \cite{draper2003} \\
FGF & \emph{fgf3} & ligand & in situ: dorsolateral margin from late blastula to early gastrula; overexpression expands the neural region and suppresses anterior neural fate, with little mesoderm-inducing activity & strong & FGF & \cite{koshida2002} \\
\rowcolor{gray!10}FGF & \emph{fgf4} & ligand & in situ: posterior notochord and Kupffer's vesicle region in early segmentation; knockdown randomizes visceral left-right patterning, removes lefty1 in the posterior notochord and impairs Kupffer's vesicle cilia & strong & FGF & \cite{yamauchi2009} \\
RA & \emph{aldh1a2} & synthesis & in situ: involuting mesoderm during gastrulation, later paraxial mesoderm, branchial arches, eyes and fin buds; neckless and no-fin (raldh2) mutants: hindbrain/spinal cord AP defects and no pectoral fin, rescued by RA; RA from paraxial mesoderm acts non-autonomously on the neural tube & strong & RA (sending) & \cite{grandel2002,begemann2001,dobbsmcauliffe2004} \\
\rowcolor{gray!10}Notch & \emph{dla} & ligand & in situ: large patches of contiguous cells in the neural plate (with deltaD) within which primary neurons are singled out; misexpression delivers lateral inhibition; itself repressed by lateral inhibition & strong & Notch (sending) & \cite{haddon1998} \\
Notch & \emph{dlb} & ligand & in situ: scattered cells within the deltaA/deltaD neural plate patches that become primary neurons; misexpression delivers lateral inhibition & strong & Notch (sending) & \cite{haddon1998} \\
\rowcolor{gray!10}Notch & \emph{dlc} & ligand & in situ: mesoderm from 50\% epiboly, tailbud, stripes in the presomitic mesoderm, posterior half of each somite, notochord; also retina, ear and scattered epidermal cells; oscillates in the PSM; beamter/deltaC mutants and morphants disrupt the segmentation clock (her1 oscillation) & strong & Notch (sending) & \cite{smithers2000,holley2002,julich2005} \\
Notch & \emph{dld} & ligand & in situ: neural plate patches (with deltaA) and somite development; misexpression delivers lateral inhibition; after eight (aei)/deltaD mutants lose her1 oscillation in the PSM; deltaD itself does not oscillate & strong & Notch (sending) & \cite{haddon1998,holley2000,holley2002} \\
\rowcolor{gray!10}Hedgehog & \emph{shha} & ligand & in situ: notochord, floor plate (ventral midline of neural tube and brain, incl. ventral diencephalon) and fin mesoderm; ectopic expression consistent with floor plate induction and promotes proximal eye fates; sonic-you mutants: defects in somite patterning, lateral floor plate, fins and motor axons (medial floor plate still forms) & strong & Hedgehog (sending) & \cite{krauss1993,ekker1995,schauerte1998} \\
Hedgehog & \emph{shhb} & ligand & in situ: ventral midline of neural tube and brain incl. floor of the diencephalon, not notochord; ectopic expression promotes proximal eye fates; lost with the ventral forebrain in cyclops mutants & strong & Hedgehog (sending) & \cite{ekker1995} \\
\midrule
\multicolumn{7}{@{}l}{\textit{secreted antagonists and degradation (blocking)}} \\
Nodal & \emph{lft1} & secreted antagonist & overexpression abolishes mesoderm induction, suppressed by Cyclops/Squint or Activin; expression depends on Nodal (feedback); midline expression overlapping cyc/sqt, later left diencephalon; lft1+lft2 knockdown: unchecked Nodal signalling, expanded mesendoderm and loss of ectoderm; lft1 alone: left-right defects & strong & Nodal & \cite{thisse1999,meno1999,bisgrove1999,feldman2002} \\
\rowcolor{gray!10}Nodal & \emph{lft2} & secreted antagonist & midline expression overlapping cyc/sqt, later left lateral plate; lft2 knockdown alone has no obvious phenotype, lft1+lft2 knockdown expands mesendoderm through unchecked Squint signalling & strong & Nodal & \cite{bisgrove1999,feldman2002} \\
BMP & \emph{chd} & secreted antagonist & in situ: organizer and transiently axial mesoderm, also paraxial mesoderm, ectoderm and brain regions; domains expand in dorsalized mutants by mid-gastrula; overexpression dorsalizes; chordino (dino) mutants ventralized & strong & BMP (blocking) & \cite{millerbertoglio1997,hammerschmidt1996,kishimoto1997,dick2000} \\
\rowcolor{gray!10}BMP & \emph{nog1} & secreted antagonist & noggin1 expressed in the fish organizer and dorsalizes on overexpression; all three zebrafish noggins antagonize ventralizing BMPs & moderate & BMP (blocking) & \cite{furthauer1999,bauer1998} \\
BMP & \emph{nog2} & secreted antagonist & overexpression antagonizes ventralizing BMPs; expression differs from noggin1/noggin3 (domain not stated in the abstract) & moderate & BMP (blocking) & \cite{furthauer1999} \\
\rowcolor{gray!10}BMP & \emph{fsta} & secreted antagonist & in situ: from mid-gastrula in anterior paraxial mesoderm (head mesoderm and first five somites), transient in the posterior notochord anlage, later brain, eyes and somites; excluded from the organizer; domain regulated by Bmp2b and Chordino; overexpression dorsalizes like Xenopus follistatin & strong & BMP (blocking) & \cite{bauer1998} \\
BMP & \emph{szl} & secreted antagonist & in situ: ventral from late blastula through gastrula, posterior end during segmentation; ogon mutants/morphants ventralized; overexpression dorsalizes, needs Chordin; does not inhibit Wnt8-dependent transcription & strong & BMP & \cite{yabe2003,martyn2003} \\
\rowcolor{gray!10}Wnt & \emph{dkk1b} & secreted antagonist & in situ: dorsal marginal blastoderm and dorsal YSL after MBT, entire margin at late blastula, embryonic shield during gastrulation, then anterior axial mesendoderm (prospective prechordal plate); maintenance reduced in boz, sqt and oep mutants; overexpression enlarges forebrain, eyes and axial mesendoderm & strong & Wnt (blocking) & \cite{hashimoto2000dkk1} \\
Wnt & \emph{frzb} & secreted antagonist & sfrp3 (frzb) cloned; expressed mainly in mesodermal and endodermal tissues (stages not given in the abstract) & moderate & Wnt (blocking) & \cite{tendeng2006} \\
\rowcolor{gray!10}Wnt & \emph{sfrp1a} & secreted antagonist & in situ: anterior neural plate including the anterior neural border cells (with tlc), complementary to posterior wnt expression; later ventral forebrain, midbrain and hindbrain; sfrp1a loss of function upregulates canonical Wnt signalling (hematopoietic stem cells) & strong & Wnt (blocking) & \cite{tendeng2006,ide2026} \\
RA & \emph{cyp26a1} & de\-gra\-da\-tion & in situ: presumptive anterior neural ectoderm and blastoderm margin from late blastula, circumblastoporal ring during gastrulation, tailbud during somitogenesis; complementary to hoxb1b; knockdown or giraffe mutation expands posterior (spinal cord) fates; encodes the RA 4-hydroxylase & strong & RA (blocking) & \cite{kudoh2002,dobbsmcauliffe2004,emoto2005,white1996,martin2010} \\
\end{longtable}
